\section{Results}
\label{sec:results}

Nine measurements, in the order the argument needs them. Each is stated here with its numbers;
the tables, the funnel columns behind them and the reproduction records are in
Appendix~\ref{app:extra}, and every scope condition is collected in \S\ref{sec:limitations}.

\subsection{BFCL v4: the score moves the whole span of the metric, and repair is complementary}
\label{sec:bfcl}

BFCL~v4 supplies official data, a multi-turn executor and a scorer at pinned Gorilla commit
\texttt{6ea57973c7a6097fd7c5915698c54c17c5b1b6c8}. Holding weights, cases, decoding, seeds,
executor and scorer fixed on 200 cases and changing \emph{only the serving adapter
configuration}, BFCL's own scorer reports \textbf{0.00} or \textbf{0.96 / 0.19}
(\texttt{simple\_python} / \texttt{multi\_turn\_base}) for the same model. The intervention is a
pair --- parser \emph{and} its chat template --- so we ran the $2\times2$ that separates them.

\begin{center}
\small
\captionof{table}{Template $\times$ parser on BFCL, first request per case.}\label{tab:bfcl2x2}
\begin{tabular}{@{}lll@{}}
\toprule
& documented template & dedicated template \\
\midrule
\texttt{hermes} parser & 0/200 & \textbf{0/200} \\
dedicated parser & \textbf{0/200} & 196/200 \\
\bottomrule
\end{tabular}
\end{center}

\textbf{Starting from the documented configuration, either one-component replacement recovers
0/200 calls, whereas replacing both components recovers 196/200.} These are simple effects at a
specified baseline, not zero marginal main effects. The $2\times2$ establishes strong
complementarity in parsing; BFCL scores were computed only for the two end-to-end corner
configurations, so the score swing itself is not decomposed over all four cells. Under the documented template the model writes a bare call inside a
\texttt{```json} fence; under the dedicated template it writes the same payload, byte for byte,
inside \texttt{<tools>} tags. Each parser behaves correctly for the envelope it is defined to
accept, so no component is defective and repairing one side of the contract buys exactly nothing.
The funnel behind those scores: the documented arm parses \textbf{0 of 200} first requests with
zero HTTP errors, while \textbf{166} of those same responses carry a call meeting the strict
criterion; the repaired arm parses 196, executes 98 of 100 \texttt{multi\_turn\_base} cases and
succeeds on \textbf{19}, against 0 executed and 0 succeeded for \texttt{hermes}, whose zero is
therefore \emph{by construction rather than by capability}
(Tables~\ref{tab:bfclfunnel}--\ref{tab:bfcl5}). \textbf{A published benchmark number for this
model can differ by the entire span of the metric depending on a serving flag the model card does
not mention.} We claim only that 0.00 is not a measurement of the model; and since the repaired
arm uses the community adapter, not a verified-optimal one, and still fails 4 of 200. We therefore
report $0\to19$ as the recovery achieved by this adapter, not as an estimate of an optimal model
or interface score.

BFCL's shipped handler for local Qwen models bypasses the serving parser entirely, so we route
BFCL's own OpenAI handler through \texttt{/v1/chat/completions}, adding only
\texttt{tool\_choice: auto} and trace headers; dataset loading, execution and scoring remain
BFCL's code and \texttt{bfcl\_eval}'s files are byte-identical (Appendix~\ref{app:bfcl}).

\subsection{$\tau$-bench: the mechanism survives a simulated user and a stateful environment}
\label{sec:tau}

BFCL scores a single trajectory against a final state. $\tau$-bench \citep{yao2025taubench} adds a
simulated user and a stateful environment, so a censored call costs not one action but the rest of
the conversation. Across all 115 \texttt{retail} tasks, paired across arms with none dropped and
changing only the serving adapter: server-parsed calls \textbf{0 versus 636}, tool executions
\textbf{0 versus 636}, and tasks reaching \emph{any} tool execution \textbf{0 versus 103}. Every
one of the documented arm's 1389 requests carried \texttt{tools}, not one was parsed, and 789
assistant turns carried a call-shaped payload the server discarded (Table~\ref{tab:tau}).
\textbf{The task-outcome difference does not reach significance and we do not claim it}: 7 versus
10 solved, three discordant pairs all in the repaired direction, exact McNemar \emph{p}=0.25, with
the documented arm's solved set a strict subset of the repaired arm's. Two scope conditions belong
with these numbers rather than in a footnote: the user simulator is a locally served
Llama-3.1-8B rather than the \texttt{gpt-4o} of the original harness, so \textbf{the absolute
scores are not comparable to the $\tau$-bench leaderboard} and only the between-arm difference is
claimed; and the arms differ in conversation length by construction, which is why every quantity
is per task and not per request.

\paragraph{The envelope layer is domain-invariant; the payload layer is not.}
$\tau$-bench also lets us ask whether the \emph{failure layer} travels across domains, because
its arguments are short scalars (\texttt{\{"order\_id": "W123"\}}) where our own task passes a
whole program as a string. Replaying vLLM's extractor over the documented arm's 789
emitted-but-unparsed assistant turns puts \textbf{789 of 789 (100\%) at the envelope layer}, with
\textbf{0} malformed payloads, \textbf{0} recoverable by lenient decoding, and --- as everywhere
else in this paper --- \textbf{0 genuine parser losses}. The emissions carry the same signature as
the code task: a bare call inside a \texttt{```json} fence, never the \texttt{<tool\_call>}
envelope. This matters in two directions. The layer that carries our headline result is
\textbf{unchanged by the domain}, and the running total of first turns in which the server
discarded a call its own extractor would have accepted stays at zero across code and non-code
alike. But the \emph{other} layer moves: Qwen2.5-Instruct loses 41\% of its unparsed items to
malformed payloads on the code task (Table~\ref{tab:layer}), almost always by leaving docstring
quotes or raw newlines unescaped inside a \texttt{code} string --- a failure mode that short
scalar arguments cannot produce. \textbf{We therefore expect the envelope-layer result to
generalise and the payload-layer proportions not to}, and we state that as a prediction rather
than a measurement, having tested one non-code domain and not a sample of them.

\subsection{Censoring against scale, and the same scale under a matched interface}
\label{sec:scale}

Our own probe, across the Qwen2.5-Coder ladder under \texttt{tool\_choice: auto} and
\texttt{hermes} --- the documented recommendation for Qwen2.5 --- n=100 per size.

\begin{center}
\small
\captionof{table}{Qwen2.5-Coder under the documented \texttt{hermes} configuration, n=100 per size.}\label{tab:scale}
\begin{tabular}{@{}llllll@{}}
\toprule
Size & Server-parsed & \textbf{tight} & strong ($\supseteq$ tight) & weak (disjoint) & Final pass \\
\midrule
1.5B & 0 & 0 & 0 & 27* & 31 \\
3B & 0 & 4 & 5 & 28 & 37 \\
7B & 0 & 21 & 22 & 1 & 53 \\
14B & 0 & 36 & 42 & 3 & 53 \\
\textbf{32B} & \textbf{0} & \textbf{80} & 100 & 0 & 68 \\
\bottomrule
\end{tabular}
\end{center}

Across a 21$\times$ range the server reports a flat zero while well-formed emitted calls rise to
80/100 ($\approx$72 after the 0.900 correction factor of \S\ref{sec:humanval}, assuming
tight-stratum precision transfers across sizes).
\textbf{Within this family, the larger the checkpoint, the larger the absolute undercount under
the same mismatched interface} --- a within-family observation across five checkpoints, not a law
relating capability to censoring. Raw-tag counts reproduce vLLM issue \#32926's baseline exactly:
\texttt{<tool\_call>} and \texttt{<tools>} appear zero times at every size. Two offline analyses
close the obvious objections (Appendix~\ref{app:reparse}): re-parsing the \emph{same stored bytes}
under four rules excludes ``your extractor is simply better than \texttt{hermes}'', and replaying
vLLM's extractor line for line to locate each loss shows \textbf{genuine parser loss is zero in
every arm}, across all 4254 de-duplicated first turns.\footnote{The 1.5B row uses the
\emph{optional} tool instruction; under the \emph{mandatory} prompt used for training
(\S\ref{sec:rl}) the same checkpoint gives weak = 66 and unparsable = 46. Coercion produces more
JSON-shaped debris, not more valid calls.}

\paragraph{A matched-interface supporting ladder through 8B.} The server column is flat zero, so the growth comes entirely from
the emitted column --- and larger checkpoints emit more well-formed JSON for reasons unrelated to
tool calling. Qwen3 supplies a ladder whose template injects tools and requires the
\texttt{<tool\_call>} envelope \texttt{hermes} accepts; the prediction (``parsed $>0$ at every
size'') was committed to the repository before running (\texttt{p4/PREREGISTRATION.md}).
Server-parsed counts are 86 / 44 / 70 / 97 at 0.6B--8B, and the silent fraction through 8B is
\textbf{1, 1, 2, 0} against the mismatched ladder's 0, 4, 21, 36, 80
(Table~\ref{tab:counterfactual}). \textbf{Across a comparable span of scale under a matched
envelope, we do not observe the growing silent fraction seen in Qwen2.5-Coder.} This is a
cross-family supporting observation, not a strict scale counterfactual: it changes family and
envelope together and does not independently measure capability growth. Two caveats we state rather than bury: the parsed column is \textbf{not
monotone}, so this shows presence, not a trend; and Qwen3 was served with
\texttt{enable\_thinking=false}, because with thinking on the model exhausts its token budget in a
\texttt{<think>} block before emitting a call --- a confound unrelated to the envelope, recorded
as a protocol deviation before the run.

A second family showing the \emph{same} undercount would be stronger, and we did not get one. Of
nine families surveyed beforehand, six do not inject OpenAI-style tools at all; of the remaining
three, Granite-3.1-8B does not attempt tool calls on this task --- 0 under \texttt{hermes}
\emph{and} 0 under its own \texttt{granite} parser. \textbf{That rescue arm is what tells us the
arm is uninformative rather than confirmatory}, and we claim nothing from it.

\subsection{Within-lineage control, and the ceiling it puts on multi-turn rescue}
\label{sec:instruct}

Changing family cannot test the scale result --- Llama parses correctly, Mistral returns 400,
DeepSeek never injects the template, so none can exhibit censoring at all.
Qwen2.5-\textbf{Instruct} can: it shares Coder's chat template, which is what produces the
false-positive template check of \S\ref{sec:families}, and differs in whether the checkpoint was
trained on tool-call tokens. On the same items, parser, \texttt{tool\_choice}, temperature and
seed, its server-parsed rate runs \textbf{1, 11, 63, 77, 88} against Coder's flat zero
(Table~\ref{tab:turn1}). \textbf{The contrast localises the mismatch to checkpoint-specific
training rather than to the shared family and template, or to parameter count alone.} We stop
short of naming tool-call-format training as \emph{the} cause: the two differ in their whole
post-training mixture, and this varies the bundle, not one element. It also answers the objection
that Instruct emits better-shaped calls around worse contents --- it solves 13--42\% of items on
turn~1 and 13--65\% by the end, nowhere near the zero that would require.

The multi-turn layer was initially void, because that host lacked \texttt{pytest}; we re-ran all
five arms with the executor working, changing nothing else and not touching the probe binary, and
the parsed column reproduced value for value. \textbf{Rescues then rise with the parsed column} ---
0, +3, +6, +14, +23 against 1, 11, 63, 77, 88, monotone together. Coder's first-turn and final
passes are identical at every size (31/31, 37/37, 53/53, 53/53, 68/68), which it would be natural
to read as multi-turn repair being useless at this scale; the Instruct ladder, on the same items
and criterion, shows it is not. We therefore state the repair result as a ceiling rather than an
absence: \textbf{how much multi-turn recovers is bounded by how much of the tool channel the
interface leaves open.}

\subsection{A 23\% failure that survived three controls and fell to the fourth}
\label{sec:llama}

Llama-3.1-8B under FC issues a tool call on 97/100 items, but on 23 the call names \emph{the
function the task asks it to implement}, with that function's own parameters as arguments; a re-run
restricted to those 23 reproduced 23/23. \emph{This also corrects a measurement of our own}:
because the probe did not verify the tool name, these 23 were counted as initiations, so
\textbf{Llama's true \texttt{run\_tests} initiation rate is 74/100, not 97/100}
(Appendix~\ref{app:errata}). Four single-variable paired controls: a richer schema explicitly
warning that the tool is not the task function left it at 22/100 (\emph{p}=1.000); a Thought
scaffold before the call raised it to \textbf{59} (\emph{p}\textless0.001, \emph{worse}), the model
supplying exactly the parameters it had just reasoned about; the official chat template moved
nothing (23$\to$22, wrong-tool discordances 1/0, exact \emph{p}=1.000); \texttt{strict: true}
on the official-template baseline took it from 22 to \textbf{0} (22/0, exact
\emph{p}=4.77$\times10^{-7}$). The earlier 0.125 and 0.0001 labels came from different outcome
comparisons and are withdrawn from the wrong-tool table. The first three rejected their alternative explanations and pointed at the
model; the fourth found the cause and withdrew that conclusion.

Adding \texttt{strict} to the official-template arm drives wrong-tool calls from 22 to 0, valid
execution from 73 to 97, turn-1 pass from 33 to 46, and final pass from 44 to 61. Thus the
single-variable final change is +17 (Table~\ref{tab:llama_arms}). The ReAct decomposition instead
uses terse FC (49) as its baseline and official-template-plus-strict (61) as the repaired endpoint,
so its broader interface-bundle contribution is \textbf{+12} (\emph{p}=0.0075), followed by
\textbf{protocol, +19}
(\emph{p}=0.0019), against +31 for both (Table~\ref{tab:decomposition}) --- \textbf{roughly two
fifths of what looked like a protocol effect was an interface effect}. With \texttt{strict} the two
protocols have identical execution mechanics (98 calls, 0 wrong-tool, 97 executed on both sides)
yet turn-1 differs 46 vs 61, so what remains is first-draft code quality under the protocol, not
tool-use mechanics. The defensible statement is narrower than either ``model defect'' or
``configuration bug'': \textbf{the model exhibits role confusion under unconstrained function
calling, and the interface constraint determines whether that confusion can be expressed as an
environment action.} We cannot show the tendency disappears once constrained, only that it stops
reaching the action space.

\subsection{Four layers, and what a pre-flight check can see}
\label{sec:taxonomy}
\label{sec:families}

The two families above fail at different layers with different remedies --- Qwen2.5-Coder at the
\emph{parser} layer, repaired by a dedicated adapter; Llama at the \emph{schema} layer, repaired by
\texttt{strict: true} --- and neither is caught by the obvious pre-flight test. Two further
families extend the taxonomy without supporting a quantitative comparison, and we report them as
such. \textbf{DeepSeek-Coder} fails at the \emph{template} layer: its chat template never injects
tools, which a template check does catch, and we found no remedy. \textbf{Mistral-7B-v0.3} fails at
the \emph{token} layer, repeating a \texttt{[TOOL\_CALLS]} marker that produces HTTP 400, and has
\textbf{no admissible 100-item FC run at all} --- all four configurations produce request errors
(parser only 2; official template 42; parser + strict 3; official template + strict 39), so vLLM's
own recommended configuration \emph{raises} the error rate from 2\% to 42\% and \texttt{strict},
which fixes Llama, does not help (Table~\ref{tab:taxonomy}).

Mistral is the one family whose failure announces itself; the other three return HTTP 200 with
\texttt{tool\_calls: []}, byte-identical to a model that declined to call the tool. \textbf{The
claim is not that all failures are silent; it is that the silent ones dominate and are the hardest
to attribute.} The natural pre-flight test --- does \texttt{apply\_chat\_template(tools=...)}
render the schema? --- correctly rejects DeepSeek-Coder but \emph{passes} Qwen2.5-Coder, which
inherits Hermes scaffolding from Qwen2.5-Instruct while never having been trained on it:
\textbf{protocol support is a per-checkpoint property and is not inferable from the chat
template.}

\subsection{Repair loop: how much comes back}
\label{sec:repair}
\label{sec:maintable}

On Qwen2.5-Coder-7B with \texttt{tool\_choice: auto} held fixed and only the adapter changed,
parse goes \textbf{0$\to$84}, items reaching a second turn \textbf{0$\to$37} and multi-turn rescues
\textbf{0$\to$9}: the mechanism is restored from literal zero (Table~\ref{tab:repair}). The two FC
arms have identical aggregate turn-1 pass, 53 vs 53. Because the evaluation adapter changes
parser, template and few-shot together, equality of this aggregate does not establish identical
first-turn generations and cannot by itself locate the failure layer. Llama's generation-time
constraint does move turn-1 pass (34$\to$46), but we use fixed-output replay and per-event tracing,
not the direction of aggregate score movement, to localize both mechanisms.

The pass-rate gain is \emph{not} significant (53$\to$62, \emph{p}=0.093). A pre-planned replication
on a random 300-item sample reproduces the direction (160/300$\to$178/300, 43 discordant against
25, \emph{p}=0.0385) --- nominally below 0.05, but this paper runs roughly a dozen McNemar tests
and the Bonferroni threshold is $\approx$0.0042, so \textbf{the repair gain still does not survive
correction and we rest no claim on it.} The protocol gap on the same 300 items
(178/300$\to$219/300, \emph{p}=3.1$\times$10$^{-6}$) clears it by three orders of magnitude.
Restricting to items both arms parsed, the residual protocol gap holds on Llama (\emph{p}=0.0021)
and is not detected on Qwen: $+8.4$\,pp, 95\% CI $[-0.5,+17.4]$, \emph{p}=0.118. \textbf{We do not
read $p>0.05$ as evidence of no effect} --- that interval admits a residual larger than the effect
we detect on Llama --- so the residual is established on one family and undetermined on the other
(Appendix~\ref{app:conditioned}).

\subsection{Parser repair restores RL experience; no held-out improvement detected}
\label{sec:rl}
\label{sec:norepair}
\label{sec:sufficiency}

The historical 1.5B result is over-determined: zero calls executed, but a replication under the
same mandatory prompt finds no well-formed call naming \texttt{run\_tests}. At 7B, by contrast, a
probe inside verl's AgentLoop records 45 complete bare-JSON calls in 115 generations, of which
\texttt{hermes} accepts none; its actor update then OOMs, so that
probe establishes the failure layer but not a learning effect. We therefore ran the missing
control: two preregistered 150-step GRPO arms at 7B under identical LoRA configuration, data,
prompt, seed, decoding, verifier and evaluation. The sole intervention is
\texttt{multi\_turn.format}: verl's stock \texttt{hermes} parser versus a registered parser that
accepts Qwen2.5-Coder's emitted format. Full rollout text, parser verdicts, executions and returned
observations were retained.

\begin{center}
\small
\captionof{table}{Formal 7B RL control. Event counts exclude held-out evaluation; outcome counts use the paired 542-item multi-turn split.}\label{tab:p3formal}
\begin{tabular}{@{}lrrrrr@{}}
\toprule
arm & tight emissions & accepted & executed/observed & final, step 0$\to$150 & rescues, 0$\to$150 \\
\midrule
broken FC & 8{,}689 & 10 & 10 & 430$\to$430 & 12$\to$12 \\
repaired FC & 17{,}608 & 16{,}912 & 16{,}844 & 430$\to$427 & 12$\to$12 \\
\bottomrule
\end{tabular}
\end{center}

The 16{,}912 accepted calls occur in 16{,}844 assistant-generation events: 40 generations contain
multiple calls and contribute 68 calls beyond the first. The AgentLoop dispatches one event per
generation; all 16{,}844 dispatches execute and return an observation. The tight classifier and
parser acceptance are not nested sets, so their marginal totals should not be read as a sequential
drop of 696 calls; Appendix~\ref{app:p3curve} gives the event-level cross-tabulation.

The broken channel is therefore \emph{nearly}, not literally, closed: only 10 of 8{,}689
tight-classified call emissions become actions. Parser repair changes the sampled experience by
more than three orders
of magnitude and returns an observation for every executed call, while holding the learning
problem fixed. This is direct causal evidence that the interface determines which emitted
behaviour enters RL experience; it supersedes the earlier, confounded ReAct comparison for that
claim.

It does not yield an outcome gain. Repaired FC ends at 415/542 turn-1 passes and 427/542 final
passes, versus 418/542 and 430/542 for broken FC; all three missing final passes are discordant in
the repaired-to-broken direction (paired difference $-0.55$ percentage points; approximate paired
95\% CI $[-1.18,0.07]$; exact McNemar $p=0.25$). Relative to its own step 0, repaired FC
has one gain and four losses ($p=0.375$), and multi-turn rescues remain exactly 12. We do not read
either non-significant difference as harm. \textbf{At this model scale, 150-step budget and single
seed, repairing the parser restores tool-mediated experience, but we detect no improvement in
held-out success or rescue counts.} The earlier 75-step ReAct arm points in the
same direction but is now only supporting evidence, because it changes protocol as well as
interface.

Both arms use LoRA rank 32 because 7B full-parameter RL does not fit a single 80\,GB card. The
held-out evaluation is identical across arms and bypasses both training parsers: verl's native
rollout RPC queries the current LoRA at temperature 0 for up to four 1,024-token
programmatic-feedback turns on the same 542 EvalPlus items, then appends sandboxed test output as
a user message. Each request attests the loaded global weight step, explaining why the
broken-training arm can still exhibit 12 evaluation rescues. The repaired process was interrupted
after step 116 and resumed from the step-90 checkpoint; before
resumption we verified model, optimizer, scheduler, RNG and dataloader state, source and
configuration hashes, and native checkpoint weight steps. The repeated step-90 evaluation differs
by two final-pass items (428 versus 426), which we report as test--retest variation; the formal
curve uses the original step-90 result and the weight-attested resumed step-120/150 results.
Mechanism totals combine original steps 1--90 with resumed steps 91--150 and exclude the abandoned
26-update step-91--116 branch. Appendix~\ref{app:p3curve} reports every checkpoint, training-log
diagnostics, checkpoint tensor deltas, and the machine-verifiable final-lineage manifest.
